\clearpage
\fancyhead[L]{\small CONNAISSANCE DES TEMS.}
\hypertarget{native-497-lower}{}\label{native-497-lower}
\begin{center}
\small Göttingische gelehrte Anzeigen. Stück 56. 57. Seite 553..562. 1808. April 7.\\[3pt]
\rule{.72\linewidth}{.3pt}
\end{center}

\emph{Connoissance des tems, ou des mouvemens célestes à l'usage des astronomes et des navigateurs pour l'an} 1808, publiée par le bureau des longitudes. Paris. De l'imprimerie impériale. 1806. 502 Seiten in Octav. --- Dieselbe für das Jahr 1809. Ebendas. 1807. 502 Seiten in Octav.

Mit diesen Jahrgängen der mit Recht so sehr geschätzten astronomischen Ephemeriden ist man wieder zu der gewöhnlichen Zeitrechnung zurück gekommen, nachdem 12 Jahrgänge (vom J. IV. bis XII.) in der Form des nun vergessenen republikanischen Kalenders erschienen waren. Die auch hier unverändert beibehaltene Einrichtung des astronomischen Kalenders selbst, und der damit verbundenen stehenden Artikel ist zu bekannt, als dass es nöthig wäre, uns dabei aufzuhalten. Nur in Ansehung der Tafel für die geographische Lage der vornehmsten Oerter der Erde --- die gewöhnlich von Zeit zu Zeit einige Zusätze und Verbesserungen erhält, obwohl noch nicht überall auch ganz bekannte Bestimmun%
\SourceLeaf{498}%
gen benutzt sind, und manche grobe Druckfehler von einem Jahrgange zum andern fortgepflanzt werden --- bemerken wir die mit dem J. 1809 gemachte Veränderung, dass dieselbe nicht mehr, wie bisher, nach den Ländern geordnet ist, sondern für die ganze Erde alphabetisch in Einem fortläuft. Rec. würde hier der bisherigen Einrichtung vor dieser Neuerung den Vorzug geben, da es bei manchen Gelegenheiten bequem ist, die von einem Lande vorhandenen guten Ortsbestimmungen gleich beisammen zu haben. Wir wenden uns also sogleich zu den \emph{Zusätzen}, die immer aus einer grossen Anzahl schätzbarer Abhandlungen, Beobachtungen und interessanter Notizen bestehen, und den Ephemeriden einen bleibenden Werth für den Astronomen geben.

Die Redaction der Zusätze, welche bisher von \textsc{Lalande} besorgt war, hat vom J. 1808 an \textsc{Delambre} übernommen. Man hat zugleich mit diesem Jahre angefangen, und wird künftig fortfahren, sie mit den auf der Pariser kaiserl. Sternwarte von \textsc{Bouvard} angestellten Beobachtungen zu eröffnen. Diese regelmässige Mittheilung der eben so zahlreichen als schätzbaren Beobachtungen, in derselben Manier, wie die Greenwicher bekannt gemacht werden, muss den Astronomen sehr willkommen sein, und erhöhet den Werth der Connoissance des tems nicht wenig. Im Jahrgange 1808 nehmen die Beobachtungen vom 24. September 1803 bis 21. September 1804, nach einer kurzen Beschreibung der Instrumente, 83 Seiten ein. Das vortreffliche Passage-Instrument von \textsc{Berge} hat $2\frac{1}{2}$ Meter (92 Pariser Zoll) Brennweite, und 11 Centimeter (4 Zoll) Oeffnung, und ist im September 1803 aufgestellt. Die Uhr ist von \textsc{Louis Berthoud}; der achtfussige, mit achromatischem Fernrohre versehene, Mauerquadrant von \textsc{Bird} ist derselbe, der ehemals \textsc{Lemonnier} zugehörte. Unter den Beobachtungen finden wir auch zahlreiche, bisher noch nicht benutzte, von der Ceres und Pallas. Angehängt sind Beobachtungen und parabolische Elemente des Kometen von 1804, und der beiden von 1805; von letztern jedoch nur die reducirten Längen und Breiten. An dieses astronomische Journal von \textsc{Bouvard} schliessen sich an, Beobachtungen von \textsc{Flaugerques} zu Viviers, und von \textsc{Vidal} zu Toulouse;  letztere bestehen in Culminationen des Mercur und der Venus, zum Theil in derselben Minute mit der Sonne von diesem scharfsehenden Astronomen beobachtet. Sonnenfinsternisse, Planeten- und Sternbedeckungen, von \textsc{Scarpellizi} zu Rom beobachtet, und zum Theil von \textsc{Lalande} berechnet. Von \textsc{Lalande} Beobachtungen der untern Conjunctionen der Venus von 1802 und 1804, und der Sonnenfinsterniss vom 1. Februar 1804 von mehreren Orten. Von \emph{eben demselben} ein neues Verzeichniss der eignen Bewegung von 500 Sternen; Rec. zählt darunter 28, wo die jährliche eigne Bewegung in gerader Aufsteigung, und 15, wo sie in der Abweichung eine halbe Raum-Secunde übersteigt. Allein es scheint, dass \textsc{Lalande} auf die Ungleichheit der Präcession nicht gehörig Rücksicht genommen hat, obgleich er es ausdrücklich versichert; wenigstens finden wir, dass dies beim Polarstern nicht geschehen, sondern schlechthin aus der \textsc{Lacaille}'schen Position von 1750 mit der \textsc{von Zach}'schen von 1800 die jährliche beobachtete Bewegung in gerader Aufsteigung $173''18$ geschlossen, und diese mit der für 1800 (nicht, wie es sein sollte, für 1775) berechneten Bewegung $207''52$ verglichen ist, wornach denn die jährliche eigne Bewegung ganz falsch $34''34$ gesetzt wird. --- Von \textsc{Prony} eine Abhandlung über die Berechnung der geographischen Längen und Breiten aus den Abständen vom Meridian, und Perpendikel auf dem elliptischen Sphäroid; seine Formeln sind zwar nur auf mässige Abstände anwendbar, empfehlen sich aber durch ihre Einfachheit. Von \textsc{Delambre} eine neue sinnreiche Methode, die Configuration der Jupiterstrabanten blos durch Rechnung vermittelst gewisser Hülfstafeln zu finden, wozu bisher gewöhnlich mechanische Hülfsmittel angewandt wurden. Von \emph{demselben} die Geschichte der Astronomie für 1804 und 1805. Die Entdeckung der \emph{Juno} nimmt hier, wie billig, den ersten Platz ein. Das Urtheil, dass jeder Astronom sich immer glücklich preisen würde, eine so grosse Entdeckung mit einem ganzen Jahre --- vor Entdeckung der Ceres hätte gewiss Niemand Bedenken getragen, zu sagen, mit einem ganzen Leben --- angestrengter Arbeit zu erkaufen, macht dem Französi%
\SourceLeaf{499}%
schen Astronomen Ehre. Dadurch, dass die Astronomie während sechs Jahren vier Mal so glücklich gewesen ist, eine so unschätzbare Bereicherung zu erhalten, kann wahrlich der Werth derselben im Preise nicht sinken: vielmehr müssen diese, hauptsächlich für Deutschland so glorreichen, Ereignisse den Muth und die Thätigkeit der Astronomen noch mehr anfeuern, da sie die Hoffnung auch zu künftigen ähnlichen glücklichen Erfolgen so sehr vermehren. Um so befremdender und mit jenem Urtheile gar nicht übereinstimmend ist die Kälte, mit welcher \textsc{Delambre} von dem vor sieben Jahren, noch vor Entdeckung der Ceres, in Lilienthal besonders von Herrn \textsc{von Zach} entworfenen Plane, den ganzen Himmel unter eine Anzahl Astronomen zur fortgesetzten Nachforschung zu vertheilen, sich ausdrückt, woran Theil zu nehmen auch er eingeladen war. Er hält eine solche Vergleichung für höchst beschwerlich und langweilig, wo man jede Nacht stundenlang maschinenmässig am unbeweglichen Fernrohr stehen und Sterne zählen müsste; er glaubt, Jeder, der Beruf zur Astronomie fühle, müsse dadurch ganz davon zurückgeschreckt werden; nur etwa \emph{beiläufig}, so oft man den Mond, die andern Planeten und die grössern Fixsterne beobachtet, solle man eine Begünstigung des Zufalls abwarten; nur der \emph{verdiene} eine solche Entdeckung zu machen, dem sie glücke, ohne darauf ausgegangen zu sein, während er sich mit einem grossen und \emph{wirklich nützlichen} Zweck beschäftige u. s. w. Aeusserungen dieser Art von einem Manne, der seiner vielen verdienstlichen Arbeiten wegen so grosse Achtung verdient, dürfen nicht mit Stillschweigen übergangen werden. Männer. deren Name als eine Autorität gilt, sollten sich, dünkt uns, um so mehr vor Urtheilen hüten, die im Geiste unsers engherzigen frivolen Jahrhunderts gesprochen zu sein scheinen könnten, wo man die Musen zu ehren glaubt, indem man sie blos zu Mägden unserer Bedürfnisse herabwürdigt, oder, wie einer unsrer geistreichsten Schriftsteller sich ausdrückt, wo man in der erhabensten Astronomie die Gestirne nur als Schrittzähler und Wegweiser für Pfefferflotten schätzt. Wer nicht lebendig genug von der Wissenschaft erwärmt ist, um das Grosse der Entdeckung eines bisher unbekannten, gleich der Erde unsere Sonne regelmässig umkreisenden, Weltkörpers zu fühlen, mit dessen Erscheinungen die Himmelsbeobachter noch nach Jahrtausenden sich beschäftigen werden; wessen erster Gedanke dabei ist, wozu wird eine solche Entdeckung nützen, dem haben wir weiter nichts zu sagen, als unsre Ueberzeugung, dass eben die neuen Planeten die Veranlassung eines ganz neuen Schwunges der theoretischen und physischen Astronomie sein werden. Wer aber andre Ansichten hat, und in den vom glücklichsten Erfolge gekrönten Nachforschungen unserer \textsc{Olbers} und \textsc{Harding} einen Grund zu der nicht chimärischen Hoffnung sieht, dass noch manche ähnliche glückliche Erfolge zu erwarten sind, der muss nothwendig lebhaft wünschen, dass jener Plan, welcher allein in einer von einzelnen Astronomen nicht zu umfassenden Unternehmung uns vom Zufall unabhängig machen kann, bald zur Ausführung kommen möge; dem wird es sehr erwünscht sein, zu hören, dass ein wichtiger Schritt dazu bereits geschehen ist. Mit wahrem Vergnügen geben wir die Nachricht, dass die Himmelskarten unsers Hrn. Prof. \textsc{Harding}, wodurch der Unternehmung so vortrefflich vorgearbeitet ist, der für die Wissenschaften bei uns so bedrängten Zeitumstände ungeachtet, ihrer endlichen Erscheinung sehr nahe sind. Es ist übrigens eine ganz ungegründete Vorstellung \textsc{Delambre}'s, wenn er die Arbeit bei solchen Nachforschungen für so beschwerlich, zeitraubend und langweilig hält. Werden sie nur auf eine zweckmässige Art angestellt, nicht mit feststehenden Instrumenten, sondern mit guten lichtstarken Nachtfernröhren, so kann man in kurzer Zeit, so bald einmal vollständige Karten entworfen sind, eine beträchtliche Strecke durchmustern, und wird gewiss immer in dieser Beschäftigung eine angenehme Erholung von andern mechanischen Arbeiten der beobachtenden und rechnenden Astronomie finden, auch ist es keinesweges nothwendig, diese Durchmusterung jede Nacht vollständig zu wiederholen.

Nach dieser Abschweifung, die uns ein Wort zu seiner Zeit schien, kommen wir auf die übrigen Artikel jenes Aufsatzes zurück. Ueber die Cometen von 1804 und 1805; \textsc{Schröter}'s und \textsc{Harding}'s
\SourceLeaf{500} Beobachtungen des Saturnringes; \textsc{Schröter}'s und \textsc{Herschel}'s Bestimmungen der Durchmesser der neuen Planeten; \textsc{Piazzi}'s Beobachtungen der Parallaxe einiger Fixsterne. Einige aus den Berliner astronomischen Jahrbüchern entlehnte Formeln von \textsc{Camerer} und \textsc{Olbers}. Ueber \textsc{Bowditch}'s und \textsc{Mendoza}'s Verfahren, scheinbare Mondabstände auf die wahren zu reduciren. Nachrichten aus dem zweiten Bande der \emph{Asiatic Researches} über die Methoden der Hindus, die Sinustafeln zu berechnen. Die übrigen kleinern Artikel dieses Jahrganges enthalten Anzeigen verschiedener neuer astronomischer Werke; die officiellen Verhandlungen wegen Abschaffung des republikanischen, und Wiedereinführung des Gregorianischen Kalenders; Beobachtungen der Sonnenfinsterniss vom 16.\ Juni 1806, und des letzten Merkurdurchganges; das Programm wegen des auf die Berechnung der Störungen der Pallas gesetzten Preises; einige Verbesserungen der Tafeln bei \textsc{Lalande}'s Astronomie, und einen Auszug aus den meteorologischen Beobachtungen zu Paris vom Jahr XII.

In dem \emph{Jahrgange} für 1809 nehmen die auf der kaiserlichen Sternwarte vom 23.\ Sept.\ 1804 bis zu Ende des Jahres 1805 angestellten Beobachtungen 108 S.\ ein. Die neuen Planeten, Ceres, Pallas und Juno, wurden 1804 fleissig im Meridian beobachtet, letztere bis zum 24.\ December; allein ungern vermissen wir Beobachtungen der Pallas und Ceres vom Ende des Jahres 1805. Von den beiden Cometen des Jahres 1805 finden wir hier die Original-Beobachtungen. Hierauf folgen fünfjährige astronomische Beobachtungen \textsc{Messier}'s, von 1760\mbox{. . . .}1764, an welche sich die schon in frühern Bänden der Connoissance des tems mitgetheilten anschliessen.  Man findet hier eine Sonnenfinsterniss, zwei Mondfinsternisse, mehrere Bedeckungen von Fixsternen, auch eine des Mars, und eine grosse Menge Verfinsterungen der Jupiterstrabanten. Für Liebhaber der Meteorologie heben wir die Bemerkung aus, dass am 19.\ Juli 1760 zwischen vier und fünf Uhr Nachmittags das Quecksilber-Thermometer in Paris auf 32 Grad stieg. --- Astronomische Beobachtungen zu Viviers von \textsc{Flaugerques}. --- Von \textsc{Burckhardt} Einrichtung des Räderwerks zur Darstellung der Bewegungen der Planeten; einige Bemerkungen, um bei verschiedenen Gelegenheiten den Gebrauch des \textsc{Borda}ischen Kreises schärfer, allgemeiner und bequemer zu machen; Untersuchungen über die eigne Bewegung unsers Sonnensystems aus Vergleichung der beobachteten eignen Bewegungen der Fixsterne: das Resultat davon ist, dass sich darüber noch wenig Befriedigendes ausmachen lässt. Einige ältere Chinesische Beobachtungen von Solstitien, aus einem vom Pater \textsc{Gaubil} an \textsc{Delisle} ehemals eingeschickten, und auf der kaiserl.\ Bibliothek verwahrten Manuscripte. Noch von \textsc{Burckhardt}, Vorschlag einer neuen Einrichtung der Spiegel-Telescope; er fängt die vom grossen Spiegel zurückgeworfenen Strahlen in der halben Brennweite mit einem auf der Axe senkrechten Planspiegel auf, der sie auf die Oculare in dem durchbrochenen Hohlspiegel zurückschickt. Diese Einrichtung hat vor der \textsc{Newton}'schen den Vortheil, dass das ganze Teleskop bei einerlei Brennweite des Spiegels nur halb so lang wird; dagegen wird dem grossen Spiegel durch den kleinen (der im Durchmesser halb so gross sein muss) ein Viertel des in die Röhre fallenden Lichtes entzogen: der Künstler \textsc{Caroché} wird die Ausführung des Vorschlags versuchen. --- Bemerkungen über eine von \textsc{Ducum} in Bordeaux vorgeschlagene Methode, die Breite zur See durch zwei ausser dem Meridian gemessene Höhen zu bestimmen, denen der Berichterstatter auch verschiedene eigne Zusätze über dies Problem und die \textsc{Douwes}'sche Auflösung beigefügt hat. --- Von \textsc{Oltmanns} eine schätzbare Untersuchung über die Länge von Quito, nach Mondsabständen und Verfinsterungen von Jupiterstrabanten von Hrn.\ \textsc{v. Humboldt} und verschiedenen ältern Beobachtungen; sein Resultat für diese bisher so schwankend bestimmte Länge ist $5^{\mathrm u}\,24^{\mathrm m}\,20^{\mathrm s}$ westlich von Paris. --- Ueber eine Stelle in Ptolemäus Geographie, die Projectionsart des Marinus von Tyrus betreffend, worin man sehr mit Unrecht schon die Spur von Mercators Projection zu erkennen geglaubt hatte. --- Beobachtungen der Ceres, Pallas und Juno, von \textsc{Poczobut} und \textsc{Keschka} zu Wilna. --- Von \textsc{Henry} strenge Formeln für die Parallaxe der Länge und
\SourceLeaf{501} Breite. --- Eine neue Art Mikrometer, von \textsc{Prony}. Die Idee, die hierbei zum Grunde liegt, ist so einfach, als sinnreich, und kann auch bei andern Gelegenheiten mit Nutzen angewandt werden. Um nemlich dem beweglichen Faden (dem Läufer) eine sehr langsame Bewegung zu geben, ohne dazu einer Schraube mit sehr feinen Gängen zu bedürfen, die nicht leicht auszuführen ist, und bald wandelbar wird, lässt er jenen durch eine Schraube mit ziemlich grossen Gängen bewegen, allein auf beiden Seiten dieser Gänge gibt er der Schraube auf derselben Axe andre, etwas Weniges engere oder weitere, unter sich gleiche, Gänge, wozu die Muttern in der Büchse festsitzen; auf diese Art wird die wahre Bewegung des Läufers nur von dem Unterschiede der beiden Arten von Gängen abhängen. --- Auszüge aus einigen neuen astronomischen Schriften; interessant sind die Auszüge aus einer Abhandlung von \textsc{Monteiro} über die Berechnung der Finsternisse im vierten Bande der Ephemeriden von Coimbra. --- Geschichte der Astronomie für 1807: betrifft hauptsächlich die Entdeckung der Vesta; deutsche Leser werden darüber hier nichts Neues antreffen. Von dem neuesten Cometen finden wir hier noch die ersten Beobachtungen von \textsc{Thulis} von 21. und 22. September; die Declinationen nur in ganzen Minuten. --- Einige geographische Bestimmungen im Mittelländischen Meere und an dessen Küste, von \textsc{Dionisio Alcala Galeano}. --- Den Beschluss dieses Jahrganges macht hier wieder ein Auszug aus \textsc{Bouvard}'s meteorologischen Beobachtungen vom Jahr XIII.
